\documentclass[11pt]{article}
\newcommand{\ListaTitulo}{Gabarito --- Lista 01}
% Preâmbulo compartilhado — MATD48, Listas de Exercícios 2026
% Usado por Lista{NN}.tex e Gabarito{NN}.tex via % Preâmbulo compartilhado — MATD48, Listas de Exercícios 2026
% Usado por Lista{NN}.tex e Gabarito{NN}.tex via % Preâmbulo compartilhado — MATD48, Listas de Exercícios 2026
% Usado por Lista{NN}.tex e Gabarito{NN}.tex via \input{preamble}

\usepackage[utf8]{inputenc}
\usepackage[T1]{fontenc}
\usepackage[brazil]{babel}
\usepackage{fullpage}
\usepackage{amsmath,amssymb}
\usepackage{enumitem}
\usepackage{xcolor}
\usepackage{fancyhdr}
\usepackage{titlesec}
\usepackage{listings}
\usepackage{hyperref}
\usepackage{booktabs}
\usepackage{graphicx}

\definecolor{matd48azul}{HTML}{0B4F6C}
\definecolor{matd48verde}{HTML}{2E8B57}

\titleformat{\section}{\normalfont\Large\bfseries\color{matd48azul}}{\thesection}{1em}{}
\titleformat{\subsection}{\normalfont\large\bfseries\color{matd48azul}}{\thesubsection}{1em}{}

\providecommand{\ListaTitulo}{Lista de Exercícios}

\setlength{\headheight}{14pt}
\pagestyle{fancy}
\fancyhf{}
\lhead{\small MATD48 --- Planejamento de Experimentos A}
\rhead{\small \ListaTitulo}
\cfoot{\thepage}
\renewcommand{\headrulewidth}{0.4pt}

\lstset{
  basicstyle=\ttfamily\small,
  breaklines=true,
  frame=single,
  columns=fullflexible,
  backgroundcolor=\color{gray!8},
  commentstyle=\color{matd48verde},
  keywordstyle=\color{matd48azul}\bfseries,
  language=R,
  extendedchars=true,
  literate=%
    {á}{{\'a}}1 {à}{{\`a}}1 {â}{{\^a}}1 {ã}{{\~a}}1
    {é}{{\'e}}1 {ê}{{\^e}}1 {í}{{\'i}}1 {ó}{{\'o}}1
    {ô}{{\^o}}1 {õ}{{\~o}}1 {ú}{{\'u}}1 {ç}{{\c c}}1
    {Á}{{\'A}}1 {À}{{\`A}}1 {Â}{{\^A}}1 {Ã}{{\~A}}1
    {É}{{\'E}}1 {Ê}{{\^E}}1 {Í}{{\'I}}1 {Ó}{{\'O}}1
    {Ô}{{\^O}}1 {Õ}{{\~O}}1 {Ú}{{\'U}}1 {Ç}{{\c C}}1
}

\hypersetup{colorlinks=true, linkcolor=matd48azul, urlcolor=matd48azul, citecolor=matd48azul}

\newcommand{\cabecalho}[2]{%
  \begin{center}
    {\Large\bfseries\color{matd48azul} #1}\\[0.3em]
    {\large #2}\\[0.3em]
    \rule{\linewidth}{0.6pt}
  \end{center}
  \vspace{0.5em}
}

\newenvironment{questao}[1]{\vspace{1em}\noindent\textbf{Questão #1.}\hspace{0.4em}}{\par}
\newenvironment{solucao}{\vspace{0.3em}\noindent\textit{Solução.}\hspace{0.4em}\color{black!85}}{\par\vspace{0.5em}}


\usepackage[utf8]{inputenc}
\usepackage[T1]{fontenc}
\usepackage[brazil]{babel}
\usepackage{fullpage}
\usepackage{amsmath,amssymb}
\usepackage{enumitem}
\usepackage{xcolor}
\usepackage{fancyhdr}
\usepackage{titlesec}
\usepackage{listings}
\usepackage{hyperref}
\usepackage{booktabs}
\usepackage{graphicx}

\definecolor{matd48azul}{HTML}{0B4F6C}
\definecolor{matd48verde}{HTML}{2E8B57}

\titleformat{\section}{\normalfont\Large\bfseries\color{matd48azul}}{\thesection}{1em}{}
\titleformat{\subsection}{\normalfont\large\bfseries\color{matd48azul}}{\thesubsection}{1em}{}

\providecommand{\ListaTitulo}{Lista de Exercícios}

\setlength{\headheight}{14pt}
\pagestyle{fancy}
\fancyhf{}
\lhead{\small MATD48 --- Planejamento de Experimentos A}
\rhead{\small \ListaTitulo}
\cfoot{\thepage}
\renewcommand{\headrulewidth}{0.4pt}

\lstset{
  basicstyle=\ttfamily\small,
  breaklines=true,
  frame=single,
  columns=fullflexible,
  backgroundcolor=\color{gray!8},
  commentstyle=\color{matd48verde},
  keywordstyle=\color{matd48azul}\bfseries,
  language=R,
  extendedchars=true,
  literate=%
    {á}{{\'a}}1 {à}{{\`a}}1 {â}{{\^a}}1 {ã}{{\~a}}1
    {é}{{\'e}}1 {ê}{{\^e}}1 {í}{{\'i}}1 {ó}{{\'o}}1
    {ô}{{\^o}}1 {õ}{{\~o}}1 {ú}{{\'u}}1 {ç}{{\c c}}1
    {Á}{{\'A}}1 {À}{{\`A}}1 {Â}{{\^A}}1 {Ã}{{\~A}}1
    {É}{{\'E}}1 {Ê}{{\^E}}1 {Í}{{\'I}}1 {Ó}{{\'O}}1
    {Ô}{{\^O}}1 {Õ}{{\~O}}1 {Ú}{{\'U}}1 {Ç}{{\c C}}1
}

\hypersetup{colorlinks=true, linkcolor=matd48azul, urlcolor=matd48azul, citecolor=matd48azul}

\newcommand{\cabecalho}[2]{%
  \begin{center}
    {\Large\bfseries\color{matd48azul} #1}\\[0.3em]
    {\large #2}\\[0.3em]
    \rule{\linewidth}{0.6pt}
  \end{center}
  \vspace{0.5em}
}

\newenvironment{questao}[1]{\vspace{1em}\noindent\textbf{Questão #1.}\hspace{0.4em}}{\par}
\newenvironment{solucao}{\vspace{0.3em}\noindent\textit{Solução.}\hspace{0.4em}\color{black!85}}{\par\vspace{0.5em}}


\usepackage[utf8]{inputenc}
\usepackage[T1]{fontenc}
\usepackage[brazil]{babel}
\usepackage{fullpage}
\usepackage{amsmath,amssymb}
\usepackage{enumitem}
\usepackage{xcolor}
\usepackage{fancyhdr}
\usepackage{titlesec}
\usepackage{listings}
\usepackage{hyperref}
\usepackage{booktabs}
\usepackage{graphicx}

\definecolor{matd48azul}{HTML}{0B4F6C}
\definecolor{matd48verde}{HTML}{2E8B57}

\titleformat{\section}{\normalfont\Large\bfseries\color{matd48azul}}{\thesection}{1em}{}
\titleformat{\subsection}{\normalfont\large\bfseries\color{matd48azul}}{\thesubsection}{1em}{}

\providecommand{\ListaTitulo}{Lista de Exercícios}

\setlength{\headheight}{14pt}
\pagestyle{fancy}
\fancyhf{}
\lhead{\small MATD48 --- Planejamento de Experimentos A}
\rhead{\small \ListaTitulo}
\cfoot{\thepage}
\renewcommand{\headrulewidth}{0.4pt}

\lstset{
  basicstyle=\ttfamily\small,
  breaklines=true,
  frame=single,
  columns=fullflexible,
  backgroundcolor=\color{gray!8},
  commentstyle=\color{matd48verde},
  keywordstyle=\color{matd48azul}\bfseries,
  language=R,
  extendedchars=true,
  literate=%
    {á}{{\'a}}1 {à}{{\`a}}1 {â}{{\^a}}1 {ã}{{\~a}}1
    {é}{{\'e}}1 {ê}{{\^e}}1 {í}{{\'i}}1 {ó}{{\'o}}1
    {ô}{{\^o}}1 {õ}{{\~o}}1 {ú}{{\'u}}1 {ç}{{\c c}}1
    {Á}{{\'A}}1 {À}{{\`A}}1 {Â}{{\^A}}1 {Ã}{{\~A}}1
    {É}{{\'E}}1 {Ê}{{\^E}}1 {Í}{{\'I}}1 {Ó}{{\'O}}1
    {Ô}{{\^O}}1 {Õ}{{\~O}}1 {Ú}{{\'U}}1 {Ç}{{\c C}}1
}

\hypersetup{colorlinks=true, linkcolor=matd48azul, urlcolor=matd48azul, citecolor=matd48azul}

\newcommand{\cabecalho}[2]{%
  \begin{center}
    {\Large\bfseries\color{matd48azul} #1}\\[0.3em]
    {\large #2}\\[0.3em]
    \rule{\linewidth}{0.6pt}
  \end{center}
  \vspace{0.5em}
}

\newenvironment{questao}[1]{\vspace{1em}\noindent\textbf{Questão #1.}\hspace{0.4em}}{\par}
\newenvironment{solucao}{\vspace{0.3em}\noindent\textit{Solução.}\hspace{0.4em}\color{black!85}}{\par\vspace{0.5em}}


\begin{document}

\cabecalho{Gabarito --- Lista de Exercícios 01}{Princípios do delineamento de experimentos (Capítulo 1 / Aula 01)}

\begin{questao}{1} \textbf{(Psicologia --- identificação de unidades)}
\end{questao}
\begin{solucao}
\begin{enumerate}[label=(\alph*)]
  \item A \textbf{unidade experimental} é o participante: é nele que o programa de meditação
    (o tratamento) é sorteado e aplicado. A \textbf{unidade amostral} é cada aplicação do
    questionário — há três medições (submuestras) por participante.
  \item Há \textbf{8 réplicas legítimas por programa} (8 participantes por grupo). As três
    respostas de um mesmo participante são submuestras, não réplicas independentes: elas
    compartilham a mesma unidade experimental e, portanto, tendem a ser correlacionadas entre
    si (um participante mais ansioso tende a responder "mais ansioso" nas três aplicações).
  \item Ele comete \textbf{pseudorreplicação}: ao tratar 72 observações correlacionadas como se
    fossem 72 réplicas independentes, infla artificialmente o tamanho amostral efetivo, reduz o
    erro-padrão das médias estimadas de forma ilegítima e, com isso, aumenta a taxa de falsos
    positivos (rejeita $H_0$ com mais frequência do que o nível de significância nominal
    permitiria). O procedimento correto — formalizado no Capítulo 3 (DCA com submuestreo) — é
    primeiro resumir as submuestras por unidade experimental (por exemplo, pela média) e só então
    comparar os 24 valores resultantes (8 por grupo).
\end{enumerate}
\end{solucao}

\begin{questao}{2} \textbf{(Agricultura --- fator, nível, tratamento)}
\end{questao}
\begin{solucao}
\begin{enumerate}[label=(\alph*)]
  \item Dois fatores: \textbf{Variedade} (níveis: V1, V2, V3) e \textbf{Densidade de plantio}
    (níveis: baixa, alta).
  \item $3 \times 2 = \textbf{6 tratamentos}$: (V1,baixa), (V1,alta), (V2,baixa), (V2,alta),
    (V3,baixa), (V3,alta).
  \item É um \textbf{desenho fatorial}, porque todas as combinações dos níveis dos dois fatores
    são usadas simultaneamente (arranjo cruzado), em vez de estudar cada fator isoladamente em
    experimentos separados. Isso permite estimar não só os efeitos principais da variedade e da
    densidade, mas também se o efeito de uma depende do nível da outra (\emph{interação}) — tema
    dos Capítulos 5 e 6.
\end{enumerate}
\end{solucao}

\begin{questao}{3} \textbf{(Ciência de dados --- os três pilares)}
\end{questao}
\begin{solucao}
\begin{enumerate}[label=(\alph*)]
  \item \textbf{Não.} A atribuição ao "novo algoritmo" foi determinada por uma escolha do
    próprio usuário (atualizar ou não o aplicativo), não por sorteio do pesquisador.
  \item Por exemplo: (i) usuários que atualizam o aplicativo tendem a ser mais engajados ou mais
    familiarizados com tecnologia, o que por si só pode aumentar o tempo de sessão,
    independentemente do algoritmo; (ii) a atualização pode ter trazido outras mudanças
    simultâneas (correções de bugs, nova interface) que também afetam o tempo de sessão,
    confundindo o efeito específico do algoritmo de recomendação.
  \item Desenho sugerido: \textbf{Fator} = algoritmo de recomendação (2 níveis: antigo, novo);
    \textbf{unidade experimental} = usuário; \textbf{aleatorização} = cada usuário elegível é
    sorteado, de forma independente, para receber o algoritmo antigo ou o novo (um teste A/B
    clássico), mantendo os dois grupos idênticos em todas as outras características do
    aplicativo durante o período do teste; \textbf{resposta} = tempo de sessão médio no período.
\end{enumerate}
\end{solucao}

\begin{questao}{4} \textbf{(Fontes de variação)}
\end{questao}
\begin{solucao}
\begin{enumerate}[label=(\alph*)]
  \item \textbf{(i) Efeito do tratamento} — é exatamente a comparação que o experimento foi
    desenhado para detectar.
  \item \textbf{(ii) Variação sistemática entre unidades experimentais} — existia antes de
    qualquer tratamento ser aplicado, então não pode ser efeito da ração.
  \item \textbf{(iii) Erro experimental} — variação residual entre unidades que receberam o
    \emph{mesmo} tratamento e tinham a \emph{mesma} origem conhecida (ninhada); é o que resta
    depois de contabilizar tratamento e a fonte de variação sistemática conhecida.
\end{enumerate}
\end{solucao}

\begin{questao}{5} \textbf{(Controle local / blocagem)}
\end{questao}
\begin{solucao}
\begin{enumerate}[label=(\alph*)]
  \item Cada uma das 4 turmas é um \textbf{bloco}. Dentro de cada turma, os alunos são
    sorteados aleatoriamente entre os dois métodos (por exemplo, metade da turma para
    expositivo, metade para baseado em projetos), repetindo esse sorteio independentemente em
    cada uma das 4 turmas.
  \item Parte da variação na nota final se deve a diferenças \emph{entre turmas} que nada têm a
    ver com o método de ensino (turma mais forte, horário, professor auxiliar etc.). Ao comparar
    os métodos \emph{dentro} de cada turma e só depois combinar essas comparações, essa fonte de
    variação conhecida é removida do erro experimental antes da análise — o que sobra para
    comparar os métodos é mais "limpo", tornando a comparação mais precisa (erro-padrão menor)
    do que sortear os métodos ignorando a turma, que misturaria variação entre turmas com o
    efeito do método.
  \item \textbf{Não.} A blocagem organiza \emph{onde} a aleatorização acontece (dentro de cada
    bloco), mas não a substitui: sem sortear os métodos dentro de cada turma, ainda haveria risco
    de viés de seleção dentro de cada bloco (por exemplo, se o professor colocasse os alunos mais
    motivados no método que prefere). Blocagem e aleatorização resolvem problemas diferentes e
    são usadas em conjunto.
\end{enumerate}
\end{solucao}

\begin{questao}{6} \textbf{(Dedução --- resultados potenciais)}
\end{questao}
\begin{solucao}
\begin{enumerate}[label=(\alph*)]
  \item O conjunto de índices $\{i : Z_i = 1\}$ é, por construção do mecanismo de aleatorização,
    um subconjunto de tamanho $n_1$ sorteado uniformemente entre os $\binom{N}{n_1}$ subconjuntos
    de $\{1,\dots,N\}$ --- exatamente a definição de uma amostra aleatória simples (AAS) sem
    reposição de tamanho $n_1$ retirada dos índices $1,\dots,N$. Como $\bar Y_1^{\,obs} =
    \tfrac{1}{n_1}\sum_{i:Z_i=1} Y_i(1)$ é a média dos valores $Y_i(1)$ correspondentes a esses
    índices sorteados, e os valores $Y_1(1),\dots,Y_N(1)$ são tratados como constantes fixas (não
    aleatórias), $\bar Y_1^{\,obs}$ é, por definição, a média de uma AAS sem reposição de tamanho
    $n_1$ do conjunto finito $\{Y_1(1),\dots,Y_N(1)\}$.

  \item Um resultado clássico de amostragem (válido para qualquer população finita, qualquer que
    seja sua distribuição) diz que a média de uma AAS sem reposição é um estimador não-viesado da
    média populacional: $E[\bar Y_1^{\,obs}] = \tfrac{1}{N}\sum_{i=1}^N Y_i(1) = \bar Y(1)$. Pelo
    item (a), essa propriedade se aplica diretamente. Pelo mesmo argumento aplicado ao
    complemento --- as $n_0$ unidades com $Z_i=0$ formam a AAS complementar de tamanho $n_0$ do
    conjunto $\{Y_1(0),\dots,Y_N(0)\}$ --- temos $E[\bar Y_0^{\,obs}] = \bar Y(0)$. Por
    linearidade da esperança,
    $$
    E\big[\widehat{\text{ATE}}\big] = E\big[\bar Y_1^{\,obs}\big] - E\big[\bar Y_0^{\,obs}\big]
    = \bar Y(1) - \bar Y(0) = \text{ATE}. \qquad \blacksquare
    $$

  \item Porque a única fonte de aleatoriedade explorada na demonstração é o \emph{mecanismo de
    aleatorização} (qual subconjunto de unidades é sorteado para tratamento) --- os $2N$ valores
    $Y_i(0), Y_i(1)$ são tratados como constantes fixas do início ao fim. O argumento é
    inteiramente combinatório/de amostragem finita (a mesma lógica que garante que a média de uma
    enquete por amostragem aleatória estima sem viés a média da população, sem qualquer suposição
    sobre a forma da distribuição dos valores populacionais). É por isso que a aleatorização, e
    não a normalidade dos erros, é o que fundamenta a validade causal do experimento --- a
    normalidade entra depois, apenas para viabilizar fórmulas fechadas de teste/IC, como já visto
    na comparação entre o teste $F$ e o teste de aleatorização.
\end{enumerate}
\end{solucao}

\begin{questao}{7} \textbf{(Método científico --- falseabilidade e validade externa)}
\end{questao}
\begin{solucao}
\begin{enumerate}[label=(\alph*)]
  \item Garante apenas \textbf{validade interna}. A aleatorização impecável assegura que, dentro
    do grupo de testadores beta, a diferença observada entre "recomendação antiga" e "recomendação
    nova" é uma estimativa não-viesada do efeito causal \emph{nessa amostra} (nenhuma variável,
    medida ou não, fica sistematicamente desbalanceada entre os dois braços). Isso nada diz sobre
    se esse efeito se generaliza a usuários fora do grupo beta --- essa é uma pergunta de validade
    externa, logicamente separada, e a aleatorização por si só não a responde.
  \item Uma formulação falseável precisa apontar uma observação que a contradiria caso a alegação
    seja falsa --- por exemplo: "entre os usuários do grupo beta sorteados para a recomendação
    nova, a taxa média de conversão será pelo menos 12\% maior do que entre os sorteados para a
    recomendação antiga, com essa diferença estatisticamente significativa a um nível pré-registrado
    (ex.\ $\alpha=0{,}05$)". Se o experimento (já desenhado e executado) resultar numa diferença
    nula, negativa, ou não significativa, a alegação é refutada por essa evidência --- ao contrário
    de uma alegação vaga como "a recomendação nova é melhor", que nenhum resultado numérico
    conseguiria contradizer de forma inequívoca.
  \item Exige uma suposição de \textbf{representatividade}: que o efeito do algoritmo de
    recomendação sobre testadores beta (usuários auto-selecionados, mais engajados, motivados por
    recompensa) seja o mesmo --- ou ao menos comparável --- ao efeito sobre a população geral de
    usuários do site. Essa suposição não é garantida por nenhuma propriedade do desenho
    experimental usado (a aleatorização só protege contra confusão \emph{dentro} do grupo estudado);
    ela é uma alegação substantiva sobre o domínio, que precisaria ser justificada separadamente
    (por exemplo, comparando características observáveis dos testadores beta com as da população
    geral) ou testada diretamente, expandindo o teste para uma amostra aleatória de todos os
    usuários.
\end{enumerate}
\end{solucao}

\vspace{1em}
\noindent\textbf{Nota de apoio (questão 1c) — a pseudorreplicação em código:}
\begin{lstlisting}
# Simulação simplificada: mesma logica da Aula 01 (Aulas2026/MATD48-01.Rmd)
# 24 participantes, 3 medidas cada (submuestras) -- efeito de tratamento nulo
set.seed(1)
participante <- rep(1:24, each = 3)
efeito_individual <- rep(rnorm(24, 0, 5), each = 3)
grupo <- rep(rep(c("10min","20min","30min"), each = 8), each = 3)
resposta <- 50 + efeito_individual + rnorm(72, 0, 2)  # tratamento sem efeito real

# Errado: ignora que ha so 24 UEs -> 72 "réplicas"
summary(aov(resposta ~ grupo))[[1]][["Pr(>F)"]][1]

# Certo: agrega por participante antes de comparar -> 24 UEs
media_part <- tapply(resposta, participante, mean)
grupo_part <- tapply(grupo, participante, unique)
summary(aov(media_part ~ grupo_part))[[1]][["Pr(>F)"]][1]
\end{lstlisting}
Rodando este código, o p-valor do modelo "errado" tende a ser bem menor que o do modelo
"correto" mesmo quando, por construção, não há efeito real de tratamento — evidência direta do
custo de ignorar a distinção entre unidade experimental e unidade amostral.

\end{document}
